\paragraph{Why did iterative WSP lose extrapolation accuracy?}
The prompts explicitly asked for prediction beyond the observations. In the original eight-task prompt, Astra was told, \emph{Among similarly fitting solutions, aim for exact outputs on unseen inputs across the stated domain.} The arithmetic prompt gave the equivalent instruction for \emph{equally good training fits.} Each revision repeated the original prompt and added, \emph{Seek a rule that also generalizes to examples not shown.} These instructions made extrapolation a secondary aim. The revision selector accepted weights with lower development loss, using training loss only to break a tie. That selection criterion rewards accuracy on the central interval and permits a worse continuation outside it.

The stored coefficients show the resulting tradeoff on square. In all six width-16 square networks, across both information conditions and three starts, the original reply put three or four changes of slope outside $[-3,3]$. After iteration, every change of slope lay inside $[-3,3]$. Each revised network therefore used a single affine continuation on either extrapolation interval. For the examples-only run at seed 164, the rightmost change of slope moved from 7.4 to 2.419 and the right-tail slope fell from 17.4 to 5.3915. At input 10, the reference is 100, the original network's real-arithmetic prediction is approximately 100, and the revised prediction is 46.6765. The selected development RMSE fell from 0.396 to 0.0145, while outer RMSE rose from 1.22 to 24.0. This coefficient-level analysis identifies the mechanism of that failure. The aggregate effect of the wording and selection rule would require a controlled ablation.
